\scanpage{29}
\mspage{\handwritten{- 47 -}}
% archive sheet number: 28

en particulier $\handwritten{\int_\gamma}d\handwritten{g}=0$ le long de tout \uline{circuit} dans A.

PROPOSITION 2.- \uline{Soit} $A\handwritten{\subset}K^n$ \uline{un domaine} contractile, \uline{et soit} $\handwritten{\sum_{k=1}^n f}_k\,dx_k$ \uline{une expression différentielle fermée à coefficients holomorphes dans} A. \uline{Il existe alors une fonction} $\handwritten{g}$ \uline{holomorphe dans} A \uline{et telle que} $d\handwritten{g}=\handwritten{\sum_{k=1}^n f}_k\,dx_k$.

Nous utiliserons le lemme suivant :

\uline{Lemme 1}.- \uline{Quels que soient les points} $\handwritten{a}$, $\handwritten{b}$ \uline{d'un domaine} $A\handwritten{\subset}K^n$, \uline{il existe une route a-chevée linéaire par morceaux}, \uline{contenue dans} A \uline{et ayant pour extrémités} $\handwritten{a}$ \uline{et} $\handwritten{b}$.

Il suffit de prouver que l'ensemble B des extrémités des routes linéaires par morceaux contenues dans A et d'origine A, est à la fois ouvert et fermé. \handwritten{O}r, si $\handwritten{c\in}\overline{B}$, il existe une boule S de centre $\handwritten{c}$ contenue dans A et contenant une extrémité $\handwritten{b}$ d'une route $\handwritten{\gamma}$ d'origine $\handwritten{a}$ contenue dans A. Comme le segment d'extrémités $\handwritten{b}$ et $\handwritten{c}$ est contenu dans S, donc dans A, $\handwritten{c}$ est l'extrémité d'une route linéaire par morceaux, juxta-position de $\handwritten{\gamma}$ et d'une route linéaire. De même, tout point $\handwritten{x\in}S$ est extrémité d'une route linéaire par morceaux, juxtaposition de $\handwritten{\gamma}$ et d'une route linéaire, donc $S\handwritten{\subset}C$, ce qui prouve que C est ouvert et fermé. Ce lemme étant démontré, soit $\handwritten{a}$ un point quelconque de A, $\handwritten{z}$ un point quelconque de A, $\handwritten{\gamma}$ et $\handwritten{\gamma}'$ deux routes contenues dans A d'origine $\handwritten{a}$ et d'extrémité $\handwritten{z}$ (il en existe en vertu du lemme). On a $\handwritten{\int_\gamma}(\handwritten{\sum_k f}_k\,dx_k)=\handwritten{\int_{\gamma'}}(\handwritten{\sum_k f}_k\,dx_k)$. En effet, si $\handwritten{\gamma}$ est définie dans un intervalle $\handwritten{[}b,c\handwritten{]}$, on peut, en remplaçant $\handwritten{\gamma}'$ par une route équivalente, supposer que $\handwritten{\gamma}'$ est définie dans un intervalle de la forme $\handwritten{[}c,d\handwritten{]}$ ; on a $\handwritten{\gamma}(c)=\handwritten{\gamma}'(d)=\handwritten{z}$ ; on peut donc définir $\handwritten{\gamma}''$ dans $\handwritten{[}b,d\handwritten{]}$, en posant $\handwritten{\gamma}''(t)=\handwritten{\gamma}(t)$ dans $\handwritten{[}b,c\handwritten{]}$ et $\handwritten{\gamma}''(t)=\handwritten{\gamma}'(c+d-t)$ dans $\handwritten{[}c,d\handwritten{]}$ ; comme $\handwritten{\gamma}''(d)=\handwritten{\gamma}'(c)=\inserted{\handwritten{a=}}\handwritten{\gamma}''(b)$, $\handwritten{\gamma}''$ est un \uline{circuit}, et notre assertion résulte du th.\unsure{1|2}.
